\section{The original native analytic chart maps preserve every original chart support}
\label{sec:projective-fixed-native-analytic-support}

\begin{theorem}[All original chart supports under the same native analytic comparison]
\label{thm:projective-fixed-native-analytic-support}
\lean{EnvelopeProperty.ProjectiveFixedNativeAnalyticSupport.nativeChartHomeomorph_mem_chartDomain_iff}
\leanok
\uses{thm:projective-fixed-native-analytic-values,thm:projective-fixed-native-analytic-charts,thm:projective-chart-generic-transitions}
\source{bj22-trivval:sec:tc}
Take an original native point in a nonempty original coordinate chart,
and map it to the corresponding analytic chart of the same fixed
reconstructed whole-field model. For every original coordinate index
$\ell$, its image belongs to the model's $\ell$th analytic coordinate
domain exactly when the original native scheme support belongs to the
original inverse $\ell$th coordinate open. No nonemptiness or matching
premise for this alternative chart is required.
\end{theorem}
\begin{proof}
Let the original source chart be $j$. On it take the actual original
regular section $X_\ell/X_j$. The proved native section equivalence
identifies it with the fixed model chart element
$c_{i,\ell}/c_{i,j}$, by the original generic transition identity.
The actual analytic section-value comparison identifies its value at
the native point with its value in the model affine chart. The genuine
model chart formula computes this as the normalized model value of
the $\ell$th homogeneous coordinate divided by the nonzero value of
the $j$th homogeneous coordinate. Thus it is nonzero exactly when the
model point lies in the alternative coordinate domain.

On the native scheme, an actual section has nonzero residue valuation
exactly at its genuine principal open. The actual pullback/basic-open
identity identifies the principal open of $X_\ell/X_j$ with the
intersection of the original $j$th and $\ell$th inverse coordinate
opens. The original point already belongs to the former. Its section
value is therefore nonzero exactly when its native support belongs to
the latter. Combining these two exact nonvanishing criteria proves
the claim for every original $\ell$, including those whose inverse
open is empty. Every point, section and coordinate domain is original.
\end{proof}
